\documentclass[10pt,a4paper]{amsart}
\usepackage[margin=19mm]{geometry}
\usepackage{amsmath,amssymb,mathtools}
\usepackage[scr=rsfs,cal=euler]{mathalfa}
\usepackage{xfrac}
\usepackage[unicode,hidelinks,bookmarks=false]{hyperref}
\newcommand{\ie}{i.e.\ }
\newcommand{\cf}{cf.\ }
\newcommand{\resp}{respectively\ }
\newcommand{\Subsection}{\subsection}
\providecommand{\nobreakdash}{\nobreak\hbox{-}}
\setlength{\emergencystretch}{2em}
\multlinegap0pt
\usepackage{amssymb}
\usepackage{float}
\newtheorem{propos}{Proposition}
\newtheorem{lemma}{Lemma}
\title{On components of stable connectivity of gradient-like diffeomorphisms of the 2-torus}
\author{D.A. Baranov, E.V. Nozdrinova, O.V. Pochinka}
\date{Source: arXiv:2605.17130v1 (CC BY 4.0)}
\begin{document}
\maketitle

\noindent\emph{Source and licence.} On components of stable connectivity of gradient-like diffeomorphisms of the 2-torus. Version arXiv:2605.17130v1 (CC BY 4.0), distributed by arXiv under the Creative Commons Attribution 4.0 International licence, http://creativecommons.org/licenses/by/4.0/, as recorded in the arXiv licence metadata for this exact version and shown on the abstract page https://arxiv.org/abs/2605.17130v1. Full licence notice: \texttt{licenses/arxiv-2605.17130-CC-BY-4.0.txt}.\par\smallskip

\noindent\textit{Editorial scope.} Counted unit: the article's lemma g10 (source0.tex lines 264--271). The article's complete original proof of it is delivered with it (source0.tex lines 272--302, one \texttt{proof} environment of 31 lines). No mathematics is written, repaired or re-derived: the statement and the proof are the article's own lines, in the article's own notation, and no retained line is substituted or re-flowed.  Retained uncounted as the source-local context the proof needs: lines 185--187; lines 206--213; lines 253--255.  One declared adaptation: exactly 2 ancillary strings inside the retained spans are omitted (image-inclusion commands and, where the source repeats a label, the repeated label definitions) and no other line differs from the source; the figure environments, with their placement, captions and labels, are kept, and the package records the omitted strings in fidelity receipts bound to the affected bodies.  The retained lines cite 2 works, whose entries are transcribed verbatim from the article's own bibliography source in the e-print and delivered with short editorial labels recorded in the item's reference mapping.  The retained lines contain the article's own diagram code, which the wrapper typesets with the standard tikz packages; no image file is delivered and the package declares no asset dependency.  Editorial formatting only, in addition: an AMS \texttt{amsart} preamble that restates the article's own declarations and macros used by the retained lines, namely the environments \texttt{propos}, \texttt{lemma} on separate counters, together with the standard packages those lines need.  The retained environments print in this document as Propos 1, Lemma 1 in order of appearance, whereas the article prints the same items with the numbering of its own full document.  The complete native source of the article as distributed in the arXiv e-print for this exact version is bundled unchanged as \texttt{sources/200769/source0.tex}.
		\begin{equation}\label{sum}
			\sum\limits_{q=0}^n(-1)^q C_q= \chi(M^n).
		\end{equation}

	\begin{propos}[\cite{treh}]\label{mn} Let $f:M^2 \rightarrow M^2$ be an orientation-preserving gradient-like diffeomorphism. Then there exist an orientation-preserving periodic homeomorphism $\phi_f:M^2 \rightarrow M^2$ of order $m_{f}$ and a flow $\xi^t_f:M^2 \rightarrow M^2$ equivalent to the gradient flow of a Morse-Smale function such that $\phi_f\xi^t_f=\xi^t_f\phi_f$ and $f=\phi_f\xi^1_f$. Moreover,
		\begin{enumerate}
			\item[1)] $B_{\phi_f}\subset\Omega_f=Fix_{\xi^1_f}$ and the invariant manifolds of the periodic points of $f$ coincide with the invariant manifolds of the fixed points of the flow $\xi^t_f$; 
			\item[2)] $f|_{\Omega_f}=\phi_f|_{\Omega_f}$ and $\tilde P_{\phi_f} \subset P_f \subset P_{\phi_f}$;
			\item[3)] $m_f$ is the smallest natural number $m$ such that $\Omega_f$ consists of fixed points of $f^m$ and all saddle points have positive orientation type; 
			\item[4)] the period of any saddle separatrix equals $m_f$.
		\end{enumerate}
	\end{propos}

	\begin{equation}\label{min}
		|\Omega_{g_{i}}|=\min_{g\in G_j^i}|\Omega_g|;
	\end{equation}

	\begin{lemma}\label{g10} The simplest diffeomorphism $g_{2}\in G_4^2$ has the following properties (see Fig. \ref{picg1}):
		\begin{enumerate}
			\item[1)] $\Omega_{g_{2}}$ consists of one fixed sink $\omega$, one source orbit $\alpha,g_{2}(\alpha)$ of period $2$, and one saddle orbit $\sigma,g_{2}(\sigma),g_{2}^2(\sigma)$ of period $3$;
			\item[2)] the set $K=W^u_\sigma\cup\omega$ is a non-contractible knot on the torus $\mathbb T^2$;
			\item[3)] the knots $K,g_{2}(K), g_{2}^2(K)$ have homotopy types $\pm\langle 1,0\rangle$, $\pm\langle 0,-1\rangle$, $\mp\langle 1,1\rangle$. 
		\end{enumerate}
		Any two simplest diffeomorphisms in $G_4^2$ are topologically conjugate.	
	\end{lemma}

	\begin{proof} We prove each item of the lemma in sequence.
		
		1) Since $G_4^2=\bigsqcup\limits_{i=0}^2\langle 3_i,2_2,1_0\rangle$, any diffeomorphism in $G^2_4$ necessarily has in its non-wandering set a fixed sink, a source orbit of period $2$, and some orbit of period $3$. The minimality condition \eqref{min} and formula \eqref{sum} lead to the following structure of the non-wandering set of the simplest diffeomorphism $g_2\in G_4^2$:
		$$\Omega_{g_2}=\{\omega,\alpha,g_{2}(\alpha),\sigma,g_{2}(\sigma),g_{2}^2(\sigma)\},$$
		where $\omega,\alpha,\sigma$ are a sink, a source, and a saddle, respectively.
		
		2) Since $\Omega_{g_{2}}$ contains a unique sink $\omega$, we have ${\rm cl}\,W^u_\sigma\setminus W^u_\sigma=\omega$. Hence, the set $K=W^u_\sigma\cup\omega$ is a knot on $\mathbb T^2$. We show that it is non-contractible. Assuming the contrary, we obtain that $K$ bounds a 2-disk $D\subset\mathbb T^2$. By item 4) of Proposition \ref{mn}, the map $g^3_2|_{K}$ reverses orientation. On the other hand, $g^3_2$ preserves orientation, and consequently $g_2^3(D)=\mathbb T^2\setminus D$. This contradicts the fact that the torus is not the union of two disks glued along their boundary.
		
		Thus, $K$ is a non-contractible knot.
		
		3) Let the knot $K$ have homotopy type $\langle a,b\rangle$. Since $g_{{2}\star}=A_4$, the knots $K,g_{2}(K),g_{2}^2(K)$ have the following homotopy types:
		\begin{itemize}
			\item $\langle K\rangle=\langle a,b\rangle$;
			\item $\langle g_{2}(K)\rangle=\langle a,b\rangle A_4=\langle b,b-a\rangle$;
			\item $\langle g_{2}^2(K)\rangle=\langle b,b-a\rangle A_4=\langle b-a,-a\rangle$.
		\end{itemize} 
		By item 2) proved above, the knots $K,g_{2}(K)$ are non-contractible and have a unique intersection point $\omega$. According to \cite{Rol}, the determinant of the matrix $\begin{pmatrix} a & b\\ b& b-a\end{pmatrix}$ must in this case have absolute value $1$. A direct computation yields the following possible homotopy types for $\langle a,b\rangle$: $\pm\langle 1,0\rangle$, $\pm\langle 1,1\rangle$, $\pm\langle 0,1\rangle$. Moreover, if $\langle K\rangle=\langle 1,0\rangle$, then 
		$\langle g_2(K)\rangle=\langle 0,-1\rangle,$ $\langle g^2_2(K)\rangle=\langle -1,-1\rangle$, $\langle g_2^3(K)\rangle=\langle -1,0\rangle,$ $\langle g^4_2(K)\rangle=\langle 0,1\rangle$, $\langle g^5_2(K)\rangle=\langle 1,1\rangle.$ Placing the knots $K,g_{2}(K),g_{2}^2(K)$ on $\mathbb T^2$ according to their homotopy types, we obtain the phase portrait of $g_{2}\in G^2_3$ (see Fig. \ref{picg1}).
		\begin{figure}[h]		
			\centerline{}		
			\caption{\small Phase portrait of $g_{2}\in G^2_4$}\label{picg1}	
		\end{figure}
		
		It remains to show that any two simplest diffeomorphisms in $G_4^2$ are topologically conjugate. To this end, we use the tricolor graph, whose equivalence class is a complete invariant of gradient-like diffeomorphisms (see \cite{treh}). To construct the graph, we color all stable (unstable) saddle separatrices blue (red). In each region complementary to the closure of the saddle separatrices, we choose one $g_2$-invariant curve and color it green. As a result, $\mathbb T^2$ is partitioned into triangular regions by colored curves. We associate a vertex to each such region and connect two vertices by an edge of the corresponding color if the regions share a boundary of that color. The resulting graph $\Gamma_{g_{2}}$ is called the {\it tricolor graph} of $g_{2}$ (see Fig. \ref{gr1}). 
		\begin{figure}[h]		
			\centerline{}		
			\caption{\small Tricolor graph of $g_{2}\in G^2_4$}\label{gr1}	
		\end{figure}
		
		By construction, all vertices of $\Gamma_{g_{2}}$ form a single red-green cycle, and $g_{2}$ induces a permutation of the vertices consisting of a rotation of this cycle by $\frac{\pi}{3}$. Two tricolor graphs are {\it equivalent} if there exists an isomorphism between them preserving edge colors and conjugating the permutations. The equivalence class of the graph is a complete invariant of topological conjugacy of the corresponding gradient-like diffeomorphism. Moreover, for any such graphs there exists an isomorphism commuting with the rotation, which completes the proof.	
	\end{proof}

\begin{thebibliography}{99}
\bibitem[Rol]{Rol} Rolfsen D. {\it Knots and links} // American Mathematical Soc. {\bf 346} (2003).
\bibitem[treh]{treh} Grines V. Z., Kapkaeva S. Kh., Pochinka O. V. {\it Tricolor graph as a complete topological invariant for gradient-like diffeomorphisms of surfaces} // Sbornik Mathematics. {\bf 205}:10, 19--46 (2014).
\end{thebibliography}
\end{document}
